\makeatletter
\def\input@path{{./}{../../../preamble/}{book2026-2027/preamble/}}
\makeatother
\input{preamble.tex}
\usepackage{needspace}
\hypersetup{pdftitle={Grade 10 Term 2 Discrete Probability}}
\setlength{\emergencystretch}{2em}
\title{Grade 10, Term 2, Lesson 3, Criterion C8\\Expected Return, Variance and Financial Dispersion}
\author{}
\date{}
\begin{document}
\maketitle
\setcounter{tocdepth}{1}
\setcounter{secnumdepth}{2}
\phantomsection\label{toc}
\tableofcontents
\section*{Evaluated criterion}
\textbf{C8:} Calculates and interprets the expected value and variance of discrete probability distributions as measures of return and risk in economic and financial contexts.
\section{Introduction}\label{sec:t2-c8-introduction}
All data in this activity are hypothetical teaching models. Probabilities have no monetary units. Keep exact expressions where useful and round final probabilities to four decimal places.
\subsection{A return is uncertain}
A school enterprise's daily net gain can be negative, zero or positive. A probability-weighted average summarises its distribution, but does not promise what tomorrow will bring. Let $X$ be gain in USD. The expected value, written $E(X)$ or $\mu$ (read mu), is $\sum_x xp_X(x)$. The sum sign means add one product for every supported value $x$.
\textbf{Check.} Does an expected gain of 3 USD guarantee 3 USD tomorrow? \emph{Suggested answer:} No. It is a model-based long-run average, not a guaranteed outcome.


\subsection{Calculate an average in small steps}
Multiply each possible value by its probability, then add. A negative return contributes a negative product. The result has the same units as $X$. Even for an integer count, the weighted average can be non-integer and need not lie in the support.
\textbf{Check.} Can a claim count have expected value 0.6? \emph{Suggested answer:} Yes. Individual counts remain integers; their weighted average need not be.


\subsection{Measure dispersion around the mean}
Variance is $\operatorname{Var}(X)=\sum_x(x-\mu)^2p_X(x)$. Subtract the mean, square each difference and weight it by its probability. Squaring prevents positive and negative deviations cancelling. An equivalent formula is $E(X^2)-\mu^2$. Standard deviation is $\sigma_X=\sqrt{\operatorname{Var}(X)}$, the non-negative square root. If $X$ is USD, variance is USD$^2$ and standard deviation is USD.
\textbf{Check.} Can variance be negative? \emph{Suggested answer:} No. It is a weighted sum of non-negative squared deviations.


\subsection{What this says about financial risk}
Greater variance means greater dispersion around the mean under the model. It does not alone describe loss probability, extreme losses, liquidity needs or the reliability of assumed probabilities. Equal-mean alternatives can have different variances, without one being universally best.
\textbf{Check.} Does lower variance alone prove an investment is best? \emph{Suggested answer:} No. Financial goals, downside, constraints and model reliability also matter.

\section{General worked problems}
\Needspace{14\baselineskip}
\subsection{Introductory general problem: a small score distribution}
A general discrete experiment produces a score. The variable $X$ takes values $ 0, 1, 2$ score units with respective probabilities $ 0.25, 0.5, 0.25$.

\textbf{Questions/Tasks}
\begin{enumerate}
\item Calculate expected value, variance and standard deviation using a calculation table; interpret their units. 
\end{enumerate}
\subsubsection*{Complete worked solution}
\begin{center}
\begin{tabular}{llp{0.48\linewidth}}
\toprule
Symbol & Value & Meaning in this problem \\ \midrule
$x$ & $0,1,2$ & Supported scores \\
$p_X(x)$ & $0.25,0.5,0.25$ & Respective masses \\
\bottomrule
\end{tabular}
\end{center}

\par\noindent\begin{minipage}{\linewidth}
\textbf{Step 1: Define and validate.} $X$ measures the stated quantity in score units; its support is the listed values. All probabilities are non-negative and total $ 1$. This supplied PMF needs no independent-trials assumption.
\end{minipage}\par\medskip
\par\noindent\begin{minipage}{\linewidth}
\textbf{Step 2: Organise the calculation.} \begin{center}
\begin{tabular}{rrrr}
\toprule
$x$ & $p_X(x)$ & $xp_X(x)$ & $x^2p_X(x)$ \\ \midrule
0 & 0.25 & 0 & 0 \\
1 & 0.5 & 0.5 & 0.5 \\
2 & 0.25 & 0.5 & 1 \\
Total & 1 & 1 & 1.5 \\
\bottomrule
\end{tabular}
\end{center}
The first column is in score units, the third in score units and the fourth in score units$^2$. Probabilities are unitless.
\end{minipage}\par\medskip
\par\noindent\begin{minipage}{\linewidth}
\textbf{Step 3: Calculate symbolically then substitute.} \[E(X)=\sum_xxp_X(x),\quad E(X^2)=\sum_xx^2p_X(x),\quad \operatorname{Var}(X)=E(X^2)-[E(X)]^2.\]The table gives $\mu=1$ score units and $E(X^2)=1.5$ score units$^2$. Thus $\operatorname{Var}(X)=1.5-(1)^2=0.5$ score units$^2$, and $\sigma_X=\sqrt{0.5}\approx0.7071$ score units.
\end{minipage}\par\medskip
\par\noindent\begin{minipage}{\linewidth}
\textbf{Step 4: Check and interpret.} The mean 1 lies between 0 and 2; variance is non-negative. The mean is a probability-weighted average, and standard deviation measures dispersion in the original units. Directly, $\sum_x(x-1)^2p_X(x)=1(0.25)+0(0.5)+1(0.25)=0.5$, agreeing with $E(X^2)-\mu^2$.
\end{minipage}\par\medskip
\subsection{Bridge to the general model}
\[
\mu=E(X)=\sum_xxp_X(x),\qquad
\operatorname{Var}(X)=\sum_x(x-\mu)^2p_X(x)=E(X^2)-\mu^2.
\]
These formulas apply when the indicated sums are finite; all distributions used here have finite support. In $E(X^2)$, square each value \emph{before} taking the weighted average; this is different from squaring $E(X)$. Standard deviation is $\sqrt{\operatorname{Var}(X)}$.

For a transformed monetary amount $Y=aX+b$, where $a$ converts each unit of $X$ into money and $b$ is a fixed monetary amount,
\[E(Y)=aE(X)+b,\qquad \operatorname{Var}(Y)=a^2\operatorname{Var}(X),\qquad \sigma_Y=|a|\sigma_X.\]
Adding a fixed amount moves every value equally and does not change dispersion. A negative multiplier reverses the order of values, but its square keeps variance non-negative.
\Needspace{14\baselineskip}
\subsection{Advanced general problem: same mean, different spread}
Model A has values $-2,2,6$ with probabilities $0.25,0.5,0.25$. Model B is always 2. Let $Y=3X-1$ USD when $X$ follows A.

\textbf{Questions/Tasks}
\begin{enumerate}
\item Compare the means and variances; calculate the transformed mean, variance and standard deviation.
\end{enumerate}
\subsubsection*{Complete worked solution}
\par\noindent\begin{minipage}{\linewidth}
\textbf{Step 1: Define and validate.} $X$ is a score under A, support $\{-2,2,6\}$. B has support $\{2\}$ with mass 1. Both are valid PMFs. $Y$ is money, support $\{-7,5,17\}$ USD.
\end{minipage}\par\medskip
\par\noindent\begin{minipage}{\linewidth}
\textbf{Step 2: Tabulate A.} \begin{center}
\begin{tabular}{rrrr}
\toprule
$x$ & $p_X(x)$ & $xp_X(x)$ & $x^2p_X(x)$ \\ \midrule
-2 & 0.25 & -0.5 & 1 \\
2 & 0.5 & 1 & 2 \\
6 & 0.25 & 1.5 & 9 \\
Total & 1 & 2 & 12 \\
\bottomrule
\end{tabular}
\end{center}

\end{minipage}\par\medskip
\par\noindent\begin{minipage}{\linewidth}
\textbf{Step 3: Calculate and compare.} For A, $E(X)=2$, $E(X^2)=12$, so $\operatorname{Var}(X)=12-2^2=8$. For B, $E(B)=2$ and $\operatorname{Var}(B)=(2-2)^2=0$. Their means agree but their dispersions differ.
\end{minipage}\par\medskip
\par\noindent\begin{minipage}{\linewidth}
\textbf{Step 4: Transform.} $E(Y)=aE(X)+b=3(2)-1=5$ USD. $\operatorname{Var}(Y)=a^2\operatorname{Var}(X)=3^2(8)=72$ USD$^2$, $\sigma_Y=\sqrt{72}\approx8.4853$ USD.
\end{minipage}\par\medskip
\par\noindent\begin{minipage}{\linewidth}
\textbf{Step 5: Check and interpret.} Directly $(-7)(0.25)+5(0.5)+17(0.25)=5$ USD. The fixed subtraction affects the mean but not variance. Lower spread does not by itself determine a universally preferred financial choice.
\end{minipage}\par\medskip
\section{Notation reference and practice procedure}
\begin{center}
\begin{tabular}{lp{0.72\linewidth}}
\toprule
Symbol & Meaning and units \\ \midrule
$p_X(x)$ & Mass at value $x$; no units \\
$\sum_x$ & Add over supported values \\
$E(X)=\mu$ & Mean; same units as $X$ \\
$E(X^2)$ & Weighted mean of squared values \\
$\operatorname{Var}(X)$ & Mean squared deviation; squared units \\
$\sigma_X$ & Standard deviation; original units \\
$a,b$ & Multiplier and fixed addition in $Y=aX+b$ \\
\bottomrule
\end{tabular}
\end{center}
\begin{enumerate}
\item Define the random variable in words and give its units and support.
\item List parameters and check that the model assumptions match the situation.
\item Translate ordinary language into an event, paying attention to endpoints.
\item Write the symbolic formula first, then substitute without rounding intermediate values.
\item Check bounds, total probability or a complementary calculation; interpret the answer in context.
\end{enumerate}
\Needspace{19\baselineskip}
\section{Foundational contextual practice}\label{sec:t2-c8-foundational}
\Needspace{14\baselineskip}
\subsection{School-enterprise net gain}
A school enterprise records one-day net gain. The variable $X$ takes values $ -10, 10, 30$ USD with respective probabilities $ 0.25, 0.5, 0.25$.

\textbf{Questions/Tasks}
\begin{enumerate}
\item Calculate expected value, variance and standard deviation using a calculation table; interpret their units. 
\end{enumerate}
\subsubsection*{Complete worked solution}
\par\noindent\begin{minipage}{\linewidth}
\textbf{Step 1: Define and validate.} $X$ measures the stated quantity in USD; its support is the listed values. All probabilities are non-negative and total $ 1$. This supplied PMF needs no independent-trials assumption.
\end{minipage}\par\medskip
\par\noindent\begin{minipage}{\linewidth}
\textbf{Step 2: Organise the calculation.} \begin{center}
\begin{tabular}{rrrr}
\toprule
$x$ & $p_X(x)$ & $xp_X(x)$ & $x^2p_X(x)$ \\ \midrule
-10 & 0.25 & -2.5 & 25 \\
10 & 0.5 & 5 & 50 \\
30 & 0.25 & 7.5 & 225 \\
Total & 1 & 10 & 300 \\
\bottomrule
\end{tabular}
\end{center}
The first column is in USD, the third in USD and the fourth in USD$^2$. Probabilities are unitless.
\end{minipage}\par\medskip
\par\noindent\begin{minipage}{\linewidth}
\textbf{Step 3: Calculate symbolically then substitute.} \[E(X)=\sum_xxp_X(x),\quad E(X^2)=\sum_xx^2p_X(x),\quad \operatorname{Var}(X)=E(X^2)-[E(X)]^2.\]The table gives $\mu=10$ USD and $E(X^2)=300$ USD$^2$. Thus $\operatorname{Var}(X)=300-(10)^2=200$ USD$^2$, and $\sigma_X=\sqrt{200}\approx14.1421$ USD.
\end{minipage}\par\medskip
\par\noindent\begin{minipage}{\linewidth}
\textbf{Step 4: Check and interpret.} The mean 10 lies between -10 and 30; variance is non-negative. The mean is a probability-weighted average, and standard deviation measures dispersion in the original units. The chance of a loss is $P(X<0)=0.25$. An expected gain of 10 USD does not remove that possibility.
\end{minipage}\par\medskip
\secInicio[][sec:t2-c8-foundational]
\Needspace{14\baselineskip}
\subsection{Insurance claim counts}
A small portfolio records its claim count in one day. The variable $X$ takes values $ 0, 1, 2$ claims with respective probabilities $ 0.5, 0.4, 0.1$.

\textbf{Questions/Tasks}
\begin{enumerate}
\item Calculate expected value, variance and standard deviation using a calculation table; interpret their units. 
\end{enumerate}
\subsubsection*{Complete worked solution}
\par\noindent\begin{minipage}{\linewidth}
\textbf{Step 1: Define and validate.} $X$ measures the stated quantity in claims; its support is the listed values. All probabilities are non-negative and total $ 1$. This supplied PMF needs no independent-trials assumption.
\end{minipage}\par\medskip
\par\noindent\begin{minipage}{\linewidth}
\textbf{Step 2: Organise the calculation.} \begin{center}
\begin{tabular}{rrrr}
\toprule
$x$ & $p_X(x)$ & $xp_X(x)$ & $x^2p_X(x)$ \\ \midrule
0 & 0.5 & 0 & 0 \\
1 & 0.4 & 0.4 & 0.4 \\
2 & 0.1 & 0.2 & 0.4 \\
Total & 1 & 0.6 & 0.8 \\
\bottomrule
\end{tabular}
\end{center}
The first column is in claims, the third in claims and the fourth in claims$^2$. Probabilities are unitless.
\end{minipage}\par\medskip
\par\noindent\begin{minipage}{\linewidth}
\textbf{Step 3: Calculate symbolically then substitute.} \[E(X)=\sum_xxp_X(x),\quad E(X^2)=\sum_xx^2p_X(x),\quad \operatorname{Var}(X)=E(X^2)-[E(X)]^2.\]The table gives $\mu=0.6$ claims and $E(X^2)=0.8$ claims$^2$. Thus $\operatorname{Var}(X)=0.8-(0.6)^2=0.44$ claims$^2$, and $\sigma_X=\sqrt{0.44}\approx0.6633$ claims.
\end{minipage}\par\medskip
\par\noindent\begin{minipage}{\linewidth}
\textbf{Step 4: Check and interpret.} The mean 0.6 lies between 0 and 2; variance is non-negative. The mean is a probability-weighted average, and standard deviation measures dispersion in the original units. The mean 0.6 claims is not a possible observed count. It describes an average over repetitions of the daily model.
\end{minipage}\par\medskip
\secInicio[][sec:t2-c8-foundational]
\Needspace{19\baselineskip}
\section{Intermediate contextual practice}\label{sec:t2-c8-intermediate}
\Needspace{14\baselineskip}
\subsection{Sales and a fixed operating cost}
A business counts daily sales; each sale contributes 8 USD and a daily fixed operating cost is 5 USD. The variable $X$ takes values $ 0, 1, 2$ sales with respective probabilities $ 0.2, 0.5, 0.3$.

\textbf{Questions/Tasks}
\begin{enumerate}
\item Calculate expected value, variance and standard deviation using a calculation table; interpret their units. Find the moments of net gain $Y=8X-5$ USD.
\end{enumerate}
\subsubsection*{Complete worked solution}
\par\noindent\begin{minipage}{\linewidth}
\textbf{Step 1: Define and validate.} $X$ measures the stated quantity in sales; its support is the listed values. All probabilities are non-negative and total $ 1$. This supplied PMF needs no independent-trials assumption.
\end{minipage}\par\medskip
\par\noindent\begin{minipage}{\linewidth}
\textbf{Step 2: Organise the calculation.} \begin{center}
\begin{tabular}{rrrr}
\toprule
$x$ & $p_X(x)$ & $xp_X(x)$ & $x^2p_X(x)$ \\ \midrule
0 & 0.2 & 0 & 0 \\
1 & 0.5 & 0.5 & 0.5 \\
2 & 0.3 & 0.6 & 1.2 \\
Total & 1 & 1.1 & 1.7 \\
\bottomrule
\end{tabular}
\end{center}
The first column is in sales, the third in sales and the fourth in sales$^2$. Probabilities are unitless.
\end{minipage}\par\medskip
\par\noindent\begin{minipage}{\linewidth}
\textbf{Step 3: Calculate symbolically then substitute.} \[E(X)=\sum_xxp_X(x),\quad E(X^2)=\sum_xx^2p_X(x),\quad \operatorname{Var}(X)=E(X^2)-[E(X)]^2.\]The table gives $\mu=1.1$ sales and $E(X^2)=1.7$ sales$^2$. Thus $\operatorname{Var}(X)=1.7-(1.1)^2=0.49$ sales$^2$, and $\sigma_X=\sqrt{0.49}\approx0.7000$ sales.
\end{minipage}\par\medskip
\par\noindent\begin{minipage}{\linewidth}
\textbf{Step 4: Check and interpret.} The mean 1.1 lies between 0 and 2; variance is non-negative. The mean is a probability-weighted average, and standard deviation measures dispersion in the original units. $E(Y)=8E(X)-5=8(1.1)-5=3.8$ USD. $\operatorname{Var}(Y)=8^2\operatorname{Var}(X)=64(0.49)=31.36$ USD$^2$ and $\sigma_Y=8(0.7)=5.6$ USD. The monetary support is $\{-5,3,11\}$ USD.
\end{minipage}\par\medskip
\secInicio[][sec:t2-c8-intermediate]
\Needspace{14\baselineskip}
\subsection{Equal-mean investment alternatives}
Alternative A yields 0 or 20 USD, each with probability $1/2$. B yields 5 or 15 USD, each with probability $1/2$.

\textbf{Questions/Tasks}
\begin{enumerate}
\item Use calculation tables to compare mean, variance and standard deviation, and give a limited financial conclusion.
\end{enumerate}
\subsubsection*{Complete worked solution}
\par\noindent\begin{minipage}{\linewidth}
\textbf{Step 1: Define and validate.} $A$ and $B$ are monetary gains with supports $\{0,20\}$ and $\{5,15\}$ USD. Each PMF totals 1. No joint independence assumption is needed to compare their separate distributions.
\end{minipage}\par\medskip
\par\noindent\begin{minipage}{\linewidth}
\textbf{Step 2: Tabulate.} \textbf{A:}\begin{center}
\begin{tabular}{rrrr}
\toprule
$x$ & $p_X(x)$ & $xp_X(x)$ & $x^2p_X(x)$ \\ \midrule
0 & 0.5 & 0 & 0 \\
20 & 0.5 & 10 & 200 \\
Total & 1 & 10 & 200 \\
\bottomrule
\end{tabular}
\end{center}
\textbf{B:}\begin{center}
\begin{tabular}{rrrr}
\toprule
$x$ & $p_X(x)$ & $xp_X(x)$ & $x^2p_X(x)$ \\ \midrule
5 & 0.5 & 2.5 & 12.5 \\
15 & 0.5 & 7.5 & 112.5 \\
Total & 1 & 10 & 125 \\
\bottomrule
\end{tabular}
\end{center}

\end{minipage}\par\medskip
\par\noindent\begin{minipage}{\linewidth}
\textbf{Step 3: Calculate.} $E(A)=10$ USD, $\operatorname{Var}(A)=200-10^2=100$ USD$^2$, $\sigma_A=10$ USD. $E(B)=10$ USD, $\operatorname{Var}(B)=125-10^2=25$ USD$^2$, $\sigma_B=5$ USD.
\end{minipage}\par\medskip
\par\noindent\begin{minipage}{\linewidth}
\textbf{Step 4: Check and interpret.} Both means are within their supports bounds. B has the same mean and less dispersion. A allows a higher gain and a lower gain; a recommendation still depends on objectives and constraints. Neither expected return is a promise.
\end{minipage}\par\medskip
\secInicio[][sec:t2-c8-intermediate]
\Needspace{19\baselineskip}
\section{Advanced contextual practice}\label{sec:t2-c8-advanced}
\Needspace{14\baselineskip}
\subsection{Claims and a reserve balance}
A firm has reserve 100 USD and pays 30 USD per claim; $X$ counts daily claims. The variable $X$ takes values $ 0, 1, 2$ claims with respective probabilities $ 0.6, 0.3, 0.1$.

\textbf{Questions/Tasks}
\begin{enumerate}
\item Calculate expected value, variance and standard deviation using a calculation table; interpret their units. Find moments of remaining reserve $Y=100-30X$ USD.
\end{enumerate}
\subsubsection*{Complete worked solution}
\par\noindent\begin{minipage}{\linewidth}
\textbf{Step 1: Define and validate.} $X$ measures the stated quantity in claims; its support is the listed values. All probabilities are non-negative and total $ 1$. This supplied PMF needs no independent-trials assumption.
\end{minipage}\par\medskip
\par\noindent\begin{minipage}{\linewidth}
\textbf{Step 2: Organise the calculation.} \begin{center}
\begin{tabular}{rrrr}
\toprule
$x$ & $p_X(x)$ & $xp_X(x)$ & $x^2p_X(x)$ \\ \midrule
0 & 0.6 & 0 & 0 \\
1 & 0.3 & 0.3 & 0.3 \\
2 & 0.1 & 0.2 & 0.4 \\
Total & 1 & 0.5 & 0.7 \\
\bottomrule
\end{tabular}
\end{center}
The first column is in claims, the third in claims and the fourth in claims$^2$. Probabilities are unitless.
\end{minipage}\par\medskip
\par\noindent\begin{minipage}{\linewidth}
\textbf{Step 3: Calculate symbolically then substitute.} \[E(X)=\sum_xxp_X(x),\quad E(X^2)=\sum_xx^2p_X(x),\quad \operatorname{Var}(X)=E(X^2)-[E(X)]^2.\]The table gives $\mu=0.5$ claims and $E(X^2)=0.7$ claims$^2$. Thus $\operatorname{Var}(X)=0.7-(0.5)^2=0.45$ claims$^2$, and $\sigma_X=\sqrt{0.45}\approx0.6708$ claims.
\end{minipage}\par\medskip
\par\noindent\begin{minipage}{\linewidth}
\textbf{Step 4: Check and interpret.} The mean 0.5 lies between 0 and 2; variance is non-negative. The mean is a probability-weighted average, and standard deviation measures dispersion in the original units. $E(Y)=100-30(0.5)=85$ USD. $\operatorname{Var}(Y)=(-30)^2(0.45)=405$ USD$^2$; $\sigma_Y=\sqrt{405}\approx20.1246$ USD. Reserve values are $\{100,70,40\}$ USD. A negative multiplier does not produce negative variance.
\end{minipage}\par\medskip
\secInicio[][sec:t2-c8-advanced]
\Needspace{14\baselineskip}
\subsection{What equal variance can hide}
Investment A yields $-10$ or 10 USD with equal probability. B yields $-20,0,20$ USD with probabilities $0.125,0.75,0.125$.

\textbf{Questions/Tasks}
\begin{enumerate}
\item Calculate both means and variances; compare loss probabilities and worst outcomes; evaluate the claim that equal mean and variance imply identical risk.
\end{enumerate}
\subsubsection*{Complete worked solution}
\par\noindent\begin{minipage}{\linewidth}
\textbf{Step 1: Define and validate.} $A$ and $B$ measure net gains in USD. Their listed supports are finite and each positive PMF totals 1.
\end{minipage}\par\medskip
\par\noindent\begin{minipage}{\linewidth}
\textbf{Step 2: Tabulate.} \textbf{A:}\begin{center}
\begin{tabular}{rrrr}
\toprule
$x$ & $p_X(x)$ & $xp_X(x)$ & $x^2p_X(x)$ \\ \midrule
-10 & 0.5 & -5 & 50 \\
10 & 0.5 & 5 & 50 \\
Total & 1 & 0 & 100 \\
\bottomrule
\end{tabular}
\end{center}
\textbf{B:}\begin{center}
\begin{tabular}{rrrr}
\toprule
$x$ & $p_X(x)$ & $xp_X(x)$ & $x^2p_X(x)$ \\ \midrule
-20 & 0.125 & -2.5 & 50 \\
0 & 0.75 & 0 & 0 \\
20 & 0.125 & 2.5 & 50 \\
Total & 1 & 0 & 100 \\
\bottomrule
\end{tabular}
\end{center}

\end{minipage}\par\medskip
\par\noindent\begin{minipage}{\linewidth}
\textbf{Step 3: Calculate.} $E(A)=E(B)=0$. For each, $\operatorname{Var}=E(X^2)-\mu^2=100-0^2=100$ USD$^2$ and $\sigma=10$ USD. Yet $P(A<0)=0.5$ and $P(B<0)=0.125$.
\end{minipage}\par\medskip
\par\noindent\begin{minipage}{\linewidth}
\textbf{Step 4: Check and interpret.} A loses more often but its worst loss is 10 USD; B loses less often but can lose 20 USD. Equal mean and variance do not imply identical downside or identical financial risk. Neither statistic alone determines a universally best choice.
\end{minipage}\par\medskip
\secInicio[][sec:t2-c8-advanced]
\section{Closing reference: key ideas and common errors}
Weight every value, including negative returns, by its probability. $E(X^2)$ is not $[E(X)]^2$. Variance uses squared units; standard deviation uses original units. A fixed addition leaves variance unchanged; scaling multiplies it by $a^2$. Expected values are averages, not guarantees. Dispersion does not fully describe financial risk.
\secInicio
\end{document}
